\begin{thebibliography}{99}
\bibitem{Watrous65} J. Watrous, \emph{The Theory of Quantum Information},
Cambridge University Press, 2018. Theorems~2.22 and~2.26
(Choi criteria), Corollary~3.40 (trace-norm contraction),
and Section~3.3, equation~(3.306) (hybrid comparison).
\bibitem{PLS66} T. Surawy-Stepney, J. Kahn, R. Kueng and M. Gu\c{t}\u{a},
Projected least-squares quantum process tomography,
\emph{Quantum} \textbf{6} (2022), 844;
arXiv:2107.01060v2, 18 October 2022.
\bibitem{Gutoski66} G. Gutoski,
On a measure of distance for quantum strategies,
\emph{J. Math. Phys.} \textbf{53} (2012), 032202;
arXiv:1008.4636v4, 14 March 2012.
\bibitem{NZGB67} S. G. Naik, M. Zartab, N. Gisin and M. Banik,
No-Go Theorem for Generic Simulation of Qubit Channels with Finite
Classical Resources, arXiv:2501.15807v2 (16 July 2025),
Theorems~2 and~5.
\bibitem{MB67} A. A. Mele and L. Bittel,
Optimal learning of quantum channels in diamond distance,
arXiv:2512.10214v3 (2026), Theorem~I.1 and Sections~III.1.1, IV.1.1.
\bibitem{OG67} A. Oufkir and F. Girardi,
Improved Lower Bounds for Learning Quantum Channels in Diamond
Distance, arXiv:2601.04180v4 (2026).
\bibitem{DKG67} R. Demkowicz-Dobrza\'nski, J. Ko\l ody\'nski
and M. Gu\c{t}\u{a}, The elusive Heisenberg limit in quantum-enhanced
metrology, \emph{Nature Communications} \textbf{3} (2012), 1063.
DOI: 10.1038/ncomms2067; arXiv:1201.3940v2, equations~(4)--(6)
and the classical-simulation argument in Methods.
\bibitem{YBCH68} Y.~Yang, G.~Bai, G.~Chiribella and M.~Hayashi,
Compression for quantum population coding,
\emph{IEEE Transactions on Information Theory} \textbf{64} (2018), no.~7,
doi:10.1109/TIT.2017.2788407;
arXiv:1701.03372v8.
\bibitem{SHW68} F.~Salek, M.~Hayashi and A.~Winter,
Usefulness of adaptive strategies in asymptotic quantum channel
 discrimination, \emph{Phys. Rev. A} \textbf{105} (2022), 022419;
arXiv:2011.06569v3 (17 March 2024).
\bibitem{CMW70} T.~Cooney, M.~Mosonyi and M.~M.~Wilde,
Strong converse exponents for a quantum channel discrimination problem
and quantum-feedback-assisted communication,
arXiv:1408.3373v3 (26 February 2016), p.~5 and Sections~4.1--4.2.
\bibitem{Acin70} A.~Ac\'in,
Statistical distinguishability between unitary operations,
\emph{Phys. Rev. Lett.} \textbf{87} (2001), 177901.
DOI: 10.1103/PhysRevLett.87.177901.
\bibitem{DFY70} R.~Duan, Y.~Feng and M.~Ying,
Entanglement is not necessary for perfect discrimination between
unitary operations, \emph{Phys. Rev. Lett.} \textbf{98} (2007), 100503.
DOI: 10.1103/PhysRevLett.98.100503.

\bibitem{DW72} S. Das and M. M. Wilde,
Quantum reading capacity: General definition and bounds,
\emph{IEEE Trans. Inform. Theory} \textbf{65} (2019), 7566--7583;
arXiv:1703.03706v3; doi:10.1109/TIT.2019.2929925.
\bibitem{WBHK72} M. M. Wilde, M. Berta, C. Hirche and E. Kaur,
Amortized channel divergence for asymptotic quantum channel discrimination,
\emph{Lett. Math. Phys.} \textbf{110} (2020), 2277--2336;
arXiv:1808.01498v2; doi:10.1007/s11005-020-01297-7.
\bibitem{WW72} X. Wang and M. M. Wilde,
Resource theory of asymmetric distinguishability for quantum channels,
\emph{Phys. Rev. Research} \textbf{1} (2019), 033169;
arXiv:1907.06306v2; doi:10.1103/PhysRevResearch.1.033169.
\bibitem{PPKKO72} Z. Pucha\l a, \L. Pawela, A. Krawiec, R. Kukulski and
M. Oszmaniec, Multiple-shot and unambiguous discrimination of von Neumann
measurements, \emph{Quantum} \textbf{5} (2021), 425;
arXiv:1810.05122v4; doi:10.22331/q-2021-04-06-425.


\bibitem{SZ74} M. Sedl\'ak and M. Ziman,
Optimal single-shot strategies for discrimination of quantum measurements,
\emph{Phys. Rev. A} \textbf{90} (2014), 052312;
arXiv:1408.0934; doi:10.1103/PhysRevA.90.052312.
Section VI, Example 4, equation (37).
\bibitem{PPKK74} Z. Pucha\l a, \L. Pawela, A. Krawiec and R. Kukulski,
Strategies for optimal single-shot discrimination of quantum measurements,
\emph{Phys. Rev. A} \textbf{98} (2018), 042103;
arXiv:1804.05856; doi:10.1103/PhysRevA.98.042103. Theorem 1.
\bibitem{FM75} J.~Fiur\'a\v sek and M.~Mi\v cuda,
Optimal two-copy discrimination of quantum measurements,
\emph{Phys. Rev. A} \textbf{80} (2009), 042312;
arXiv:0909.2940; doi:10.1103/PhysRevA.80.042312.

\end{thebibliography}
